\chapter{Implementation Source Code}\label{app:source}

\newcommand{\PyFile}[1]{{\raggedleft\color{Crimson}\ttfamily\bfseries\small#1\par}}

\begin{center}
Copyright {\textcopyright} 2021 Wei Ke and Ka-Hou Chan.
\end{center}

The type-checking and operational semantics of the pattern matching described in the thesis, as well as the test cases, are implemented in Python.
The source code listed is licensed under GPL-3, see the license file in \autoref{app:GPL-3}.

\section[Abstract syntax of all else]{Abstract syntax of this minimal oo language, except pattern matching}

\PyFile{asx.py}
\begin{dispPy}
#    (c) 2021 \BlindText{Wei Ke and Ka-Hou Chan}
#    license:          GPL-3
#
from collections import namedtuple

PrintStmt = namedtuple('PrintStmt', ['args'])
AssignStmt = namedtuple('AssignStmt', ['lexpr', 'expr'])
IfStmt = namedtuple('IfStmt', ['expr', 'then_stmt', 'else_stmt'])
WhileStmt = namedtuple('WhileStmt', ['expr', 'stmt'])
BlockStmt = namedtuple('BlockStmt', ['stmts'])

VarDecl = namedtuple('VarDecl', ['label', 'cla'])
VarEnd = namedtuple('VarEnd', ['label'])

VarExpr = namedtuple('VarExpr', ['label'])
AttrExpr = namedtuple('AttrExpr', ['expr', 'label'])
OpExpr = namedtuple('OpExpr', ['op', 'args'])
NewExpr = namedtuple('NewExpr', ['cla'])
AndExpr = namedtuple('AndExpr', ['left', 'right'])
OrExpr = namedtuple('OrExpr', ['left', 'right'])

Program = namedtuple('Program', ['block'])
#
# end of asx.py
# Id: asx.py 4842 2021-11-14 17:09:57Z \BlindText{wke@IPM.EDU.MO}
\end{dispPy}

\section[Abstract syntax of pattern matching]{Abstract syntax of this pattern matching}

\PyFile{pattern.py}
\begin{dispPy}
#    (c) 2021 \BlindText{Wei Ke and Ka-Hou Chan}
#    license:          GPL-3
#
from collections import namedtuple

LabeledPattern = namedtuple('LabeledPattern', ['label', 'base'])
PatternRef = namedtuple('PatternRef', ['label'])
ClassPattern = namedtuple('ClassPattern', ['cla', 'attrs'])
PatternConj = namedtuple('PatternConj', ['patterns'])
PatternDisj = namedtuple('PatternDisj', ['patterns'])
Case = namedtuple('Case', ['junc', 'stmt', 'extra'])
MatchStmt = namedtuple('MatchStmt', ['expr', 'cas'])

class Extra:
    def __init__(self, extra = None):
        self.extra = extra

    def put(self, extra):
        self.extra = extra

    def get(self):
        return self.extra
#
# end of pattern.py
# Id: pattern.py 4838 2021-11-14 12:02:46Z \BlindText{wke@IPM.EDU.MO}
\end{dispPy}

\section{Prety-printing of dictionaries}

\PyFile{pp.py}
\begin{dispPy}
#    (c) 2021 \BlindText{Wei Ke and Ka-Hou Chan}
#    license:          GPL-3
#
from collections import namedtuple

PrfxPP = namedtuple('PrefixedPP', ['prefix', 'pp'])

def pprint(s, name = None):
    def pp(s, indent):
        if isinstance(s, dict):
            if s:
                print('{')
                for k, v in s.items():
                    indent2 = indent+4
                    print('{}{!r}: '.format(' '*indent2, k), end = '')
                    pp(v, indent2)
                print(' '*indent+'}')
            else:
                print('{}')
        else:
            print('{!r}'.format(s))

    if type(s) is PrfxPP:
        p, t = s
        if p:
            if name:
                print('{} {}: '.format(p, name), end = '')
            else:
                print('{}: '.format(p), end = '')
        pp(t, 0)
    else:
        if name:
            print('{}: '.format(name), end = '')
        pp(s, 0)
#
# end of pp.py
# Id: pp.py 4838 2021-11-14 12:02:46Z \BlindText{wke@IPM.EDU.MO}
\end{dispPy}

\section{Bijective dictionary}

\PyFile{bidict.py}
\begin{dispPy}
#    (c) 2021 \BlindText{Wei Ke and Ka-Hou Chan}
#    license:          GPL-3
#
class BiDictError(Exception): pass
class UniqueInvError(BiDictError): pass
class InterInvError(BiDictError): pass

class bidict(dict):
    def __init__(self, *args, **kwargs):
        super().__init__(*args, **kwargs)
        self.inv = {}
        for k, v in self.items():
            self.inv.setdefault(v, []).append(k)

    def __setitem__(self, k, v):
        if k in self:
            x = self[k]
            self.inv[x].remove(k)
            if not self.inv[x]:
                del self.inv[x]
        super().__setitem__(k, v)
        self.inv.setdefault(v, []).append(k)

    def __delitem__(self, k):
        v = self[k]
        self.inv.setdefault(v, []).remove(k)
        if not self.inv[v]:
            del self.inv[v]
        super().__delitem__(k)

    def unique_inv(self, v):
        ks = self.inv[v]
        if len(ks) != 1:
            raise UniqueInvError()
        return ks[0]


def dict_union(ds):
    return {k: v for d in ds for k, v in d.items()}

def bidict_union(bds):
    return bidict(dict_union(bds))
#
# end of bidict.py
# Id: bidict.py 4838 2021-11-14 12:02:46Z \BlindText{wke@IPM.EDU.MO}
\end{dispPy}

\section{Graphs and graph algorithms}

\PyFile{graph.py}
\begin{dispPy}
#    (c) 2021 \BlindText{Wei Ke and Ka-Hou Chan}
#    license:          GPL-3
#
from collections import namedtuple
from queue import Queue

from bidict import bidict
from pp import PrfxPP

from subtype import NULL_TYPE, ValueSet, exists_ty_le_all, ty_inf, ty_sup

from pattern import LabeledPattern, PatternRef, ClassPattern

LayoutGraph = namedtuple('LayoutGraph', ['nodes', 'edges', 'root'])
PatternGraph = namedtuple('PatternGraph', ['layout', 'types'])
StateGraph = namedtuple('StateGraph', ['layout', 'types', 'values'])
Node = namedtuple('Node', ['id'])
Label = namedtuple('Label', ['id'])
Edges = namedtuple('Edges', ['labels', 'targets'])

class GraphError(Exception): pass
class Mismatch(GraphError):  pass
class NoUnion(GraphError):  pass
class NoConj(GraphError):  pass
class UndefRef(GraphError):  pass
class RedefRef(GraphError):  pass
class NoScope(GraphError):  pass
class UndefVar(GraphError):  pass
class UndefAttr(GraphError):  pass

def unzip2(s):
    if not s:
        return ([], [])
    else:
        return list(zip(*s))

def unzip4(s):
    if not s:
        return ([], [], [], [])
    else:
        return list(zip(*s))

def set_union(ss):
    return set().union(*ss)

def set_inter(ss):
    i = iter(ss)
    for s in i:
        return set(s).intersection(*i)

    return set()

def layout_graph_to_pp(g):
    ns, es, p = g
    pp = {}
    pp['root'] = '({})'.format(p.id)
    pp['nodes'] = sorted(u.id for u in ns)
    pp['edges'] = {u.id: [(la.id, es.targets[(u, la)].id) for la in es.labels[u]]
                   for u in sorted(ns)}
    return PrfxPP('LayoutGraph', pp)

class NodeFactory:
    def __init__(self):
        self._cur_id = 0

    def new_node(self):
        self._cur_id += 1
        return Node(self._cur_id)

NODE_FACTORY = NodeFactory()

def new_node():
    return NODE_FACTORY.new_node()

def init_state_graph():
    p = new_node()
    return StateGraph(LayoutGraph({p}, Edges({p: set()}, {}), p), {}, {})

def gc_layout(g):
    ns, es, r = g
    queue = Queue()
    queue.put(r)
    ns2 = set()
    es2 = Edges(dict(), dict())

    while not queue.empty():
        p = queue.get()
        if p not in ns2:
            ns2.add(p)
            las = set(es.labels[p])
            es2.labels[p] = las
            for la in las:
                q = es.targets[(p, la)]
                es2.targets[(p, la)] = q
                queue.put(q)

    return LayoutGraph(ns2, es2, r)

def gc_state(sg):
    g, tm, vm = sg
    g2 = gc_layout(g)
    tm2 = {p: t for (p, t) in tm.items() if p in g2.nodes}
    vm2 = {p: v for (p, v) in vm.items() if p in g2.nodes}
    return StateGraph(g2, tm2, vm2)

def extract_pattern(sg, p):
    (ns, es, r), tm, vm = sg
    g2 = gc_layout(LayoutGraph(ns, es, p))
    tm2 = {u: ValueSet({vm[u]}) if u in vm else tm[u] for u in g2.nodes}
    return PatternGraph(g2, tm2)

def swing_layout(g, p, la, q):
    ns, es, r = g
    es.labels[p].add(la)
    es.targets[(p, la)] = q

def swing_state(sg, p, la, q):
    swing_layout(sg.layout, p, la, q)

def add_object_to_layout(g, cla):
    ns, es, r = g
    p = new_node()
    ns.add(p)
    las = set(cla.attrs)
    es.labels[p] = las
    qs = []
    for la in las:
        q = new_node()
        ns.add(q)
        es.labels[q] = set()
        es.targets[(p, la)] = q
        qs.append(q)

    return (p, qs)

def add_object_to_state(sg, cla):
    g, tm, vm = sg
    p, qs  = add_object_to_layout(g, cla)
    tm[p] = cla
    for q in qs:
        tm[q] = NULL_TYPE

    return p

def add_value_to_state(sg, v):
    p  = add_object_to_state(sg, v.cla)
    sg.values[p] = v
    return p

def push_state(sg, sla):
    (ns, es, r), tm, vm = sg
    r2 = new_node()
    ns.add(r2)
    es.labels[r2] = {sla}
    es.targets[(r2, sla)] = r
    return StateGraph(LayoutGraph(ns, es, r2), tm, vm)

def pop_state(sg, sla):
    (ns, es, r), tm, vm = sg
    if sla not in es.labels[r]:
        raise NoScope()

    return gc_state(StateGraph(LayoutGraph(ns, es, es.targets[(r, sla)]), tm, vm))

def find_lvar(sg, sla, la):
    (ns, es, r), tm, vm = sg
    while la not in es.labels[r] and sla in es.labels[r]:
        r = es.targets[(r, sla)]

    if la not in es.labels[r]:
        raise UndefVar()

    return (r, la)

def find_var(sg, sla, la):
    return sg.layout.edges.targets[find_lvar(sg, sla, la)]

def find_lattr(sg, p, la):
    (ns, es, r), tm, vm = sg
    if not p in ns:
        raise UndefVar()

    if la not in es.labels[p]:
        raise UndefAttr()

    return (p, la)

def find_attr(sg, p, la):
    return sg.layout.edges.targets[find_lattr(sg, p, la)]

def cons_match(g1, g2, le, tau1, tau2):
    ns1, es1, p1 = g1
    ns2, es2, p2 = g2
    f = bidict()

    def dfs_match(es1, p1, es2, p2):
        if p1 in f:
            if f[p1] != p2:
                raise Mismatch()
        else:
            if not le(tau2[p2], tau1[p1]):
                raise Mismatch()
            if p2 in f.inv:
                raise Mismatch()

            f[p1] = p2

            for la in es1.labels[p1]:
                q1 = es1.targets[(p1, la)]
                if la not in es2.labels[p2]:
                    raise Mismatch()
                q2 = es2.targets[(p2, la)]
                dfs_match(es1, q1, es2, q2)

    dfs_match(es1, p1, es2, p2)
    return f

def cons_union(gs):
    ess, ps = list(zip(*((es, p) for ns, es, p in gs)))
    fs = [bidict() for g in gs]

    def dfs_union(iz, ps):
        kz = [k for k in iz if ps[k] in fs[k]]
        cz = [c for c in iz if ps[c] not in fs[c]]
        uz = [fs[k][ps[k]] for k in kz]
        if len(uz) > 1:
            raise NoUnion()
        if cz:
            if uz:
                pd = uz[0]
                if any(pd in fs[c].inv for c in cz):
                    raise NoUnion()
            else:
                pd = new_node()
            for c in cz:
                fs[c][ps[c]] = pd
            laz = set_union(ess[i].labels[ps[i]] for i in iz)
            for la in laz:
                qs = [None]*len(gs)
                jz = [j for j in iz if la in ess[j].labels[ps[j]]]
                for j in jz:
                    qs[j] = ess[j].targets[(ps[j], la)]
                dfs_union(jz, qs)

    dfs_union(list(range(len(gs))), ps)
    return fs

def cons_inter(gs):
    ess, ps = unzip2((es, p) for ns, es, p in gs)
    fs = [bidict() for g in gs]

    def dfs_inter(ps):
        uz = [f[p] for f, p in zip(fs, ps) if p in f]
        if uz:
            pc = uz[0]
        else:
            pc = new_node()
        for f, p in zip(fs, ps):
            f[p] = pc
        if len(uz) != len(ps):
            laz = set_inter(es.labels[p] for es, p in zip(ess, ps))
            for la in laz:
                qs = [es.targets[(p, la)] for es, p in zip(ess, ps)]
                dfs_inter(qs)

    dfs_inter(ps)
    return fs

def cons_match_conj(pgs):
    gs, tms = unzip2(pgs)
    fs = cons_union(gs)
    tsd = {}
    nsd = set_union(f.inv for f in fs)
    for ud in nsd:
        ts = {tm[f.unique_inv(ud)] for tm, f in zip(tms, fs) if ud in f.inv}
        if not exists_ty_le_all(ts):
            raise NoConj()
        td = ty_inf(ts)
        tsd[ud] = td

    return (fs, tsd)

def cons_match_disj(pgs):
    gs, tms = unzip2(pgs)
    fs = cons_inter(gs)
    tsc = {}
    nsc = set_union(f.inv for f in fs)
    for uc in nsc:
        ts = set_union({tm[u] for u in f.inv[uc]} for tm, f in zip(tms, fs))
        tc = ty_sup(ts)
        tsc[uc] = tc

    return (fs, tsc)

def cons_pattern_graph(pattern):
    ns = set()
    es = Edges({}, {})
    tm = {}
    rm = bidict()

    def parse(pattern):
        if type(pattern) is LabeledPattern:
            la = pattern.label
            p = parse(pattern.base)
            if la in rm:
                r = rm[la]
                if r in ns:
                    raise RedefRef()
                'E[p/r]'
                nlas = [nla for nla, u in es.targets.items() if u == r]
                for nla in nlas:
                    es.targets[nla] = p
                'R[p/r]'
                las2 = [la2 for la2, u in rm.items() if u == r]
                for la2 in las2:
                    rm[la2] = p

            rm[la] = p
        elif type(pattern) is ClassPattern:
            cla, attrs = pattern
            p = new_node()
            ns.add(p)
            tm[p] = cla
            las, qs = unzip2([(la, parse(pattern2))
                              for la, pattern2 in attrs.items()])
            es.labels[p] = set(las)
            for la, q in zip(las, qs):
                es.targets[(p, la)] = q
        elif type(pattern) is PatternRef:
            la = pattern.label
            if la in rm:
                p = rm[la]
            else:
                p = new_node()
                rm[la] = p
        elif type(pattern) is ValueSet:
            p = new_node()
            ns.add(p)
            es.labels[p] = set()
            tm[p] = pattern

        return p

    p = parse(pattern)
    if any(u not in ns for u in rm.inv):
        raise UndefRef()

    return (LayoutGraph(ns, es, p), tm, rm)
#
# end of graph.py
# Id: graph.py 4838 2021-11-14 12:02:46Z \BlindText{wke@IPM.EDU.MO}
\end{dispPy}

\section{Operations on values}

\PyFile{op.py}
\begin{dispPy}
#    (c) 2021 \BlindText{Wei Ke and Ka-Hou Chan}
#    license:          GPL-3
#
from collections import namedtuple

import operator

from subtype import Value, INT_TYPE, STR_TYPE, BOOL_TYPE

from graph import Label, add_value_to_state

OpDef = namedtuple('OpDef', ['op', 'par_types', 'res_type', 'f'])

OP_TAB = {}

def is_op(la):
    return la in OP_TAB

def get_op(la):
    return OP_TAB[la]

def invoke_op(sg, op, args):
    f = OP_TAB[op].f
    return f(sg, *args)

def op_binary(sg, x, y, t, f):
    g, tm, vm = sg
    return add_value_to_state(sg, Value(t, f(vm[x].value, vm[y].value)))

def op_unary(sg, x, t, f):
    g, ts, vs = sg
    return add_value_to_state(sg, t, Value(t, f(vs[x].value)))

def op_add(sg, x, y):
    return op_binary(sg, x, y, INT_TYPE, operator.add)

def op_sub(sg, x, y):
    return op_binary(sg, x, y, INT_TYPE, operator.sub)

def op_mul(sg, x, y):
    return op_binary(sg, x, y, INT_TYPE, operator.mul)

def op_div(sg, x, y):
    return op_binary(sg, x, y, INT_TYPE, operator.floordiv)

def op_mod(sg, x, y):
    return op_binary(sg, x, y, INT_TYPE, operator.mod)

def op_neg(sg, x):
    return op_unary(sg, x, INT_TYPE, operator.neg)

def op_not(sg, x):
    return op_unary(sg, x, BOOL_TYPE, operator.not_)

def op_cat(sg, x, y):
    return op_binary(sg, x, y, STR_TYPE, operator.concat)

def op_lower(sg, x):
    return op_unary(sg, x, STR_TYPE, lambda x: x.lower())

def op_upper(sg, x):
    return op_unary(sg, x, STR_TYPE, lambda x: x.lower())

def op_ieq(sg, x, y):
    return op_binary(sg, x, y, BOOL_TYPE, operator.eq)

def op_ine(sg, x, y):
    return op_binary(sg, x, y, BOOL_TYPE, operator.ne)

def op_ilt(sg, x, y):
    return op_binary(sg, x, y, BOOL_TYPE, operator.lt)

def op_ile(sg, x, y):
    return op_binary(sg, x, y, BOOL_TYPE, operator.le)

def op_igt(sg, x, y):
    return op_binary(sg, x, y, BOOL_TYPE, operator.gt)

def op_ige(sg, x, y):
    return op_binary(sg, x, y, BOOL_TYPE, operator.ge)

def op_seq(sg, x, y):
    return op_binary(sg, x, y, BOOL_TYPE, operator.eq)

def op_sne(sg, x, y):
    return op_binary(sg, x, y, BOOL_TYPE, operator.ne)

def op_slt(sg, x, y):
    return op_binary(sg, x, y, BOOL_TYPE, operator.lt)

def op_sle(sg, x, y):
    return op_binary(sg, x, y, BOOL_TYPE, operator.le)

def op_sgt(sg, x, y):
    return op_binary(sg, x, y, BOOL_TYPE, operator.gt)

def op_sge(sg, x, y):
    return op_binary(sg, x, y, BOOL_TYPE, operator.ge)

OP_TAB = {
    Label('add'): OpDef(Label('add'), [INT_TYPE, INT_TYPE], INT_TYPE, op_add),
    Label('sub'): OpDef(Label('sub'), [INT_TYPE, INT_TYPE], INT_TYPE, op_sub),
    Label('mul'): OpDef(Label('mul'), [INT_TYPE, INT_TYPE], INT_TYPE, op_mul),
    Label('div'): OpDef(Label('div'), [INT_TYPE, INT_TYPE], INT_TYPE, op_div),
    Label('mod'): OpDef(Label('mod'), [INT_TYPE, INT_TYPE], INT_TYPE, op_mod),
    Label('neg'): OpDef(Label('neg'), [INT_TYPE], INT_TYPE, op_neg),
    Label('not'): OpDef(Label('not'), [BOOL_TYPE], BOOL_TYPE, op_not),
    Label('cat'): OpDef(Label('cat'), [STR_TYPE, STR_TYPE], STR_TYPE, op_cat),
    Label('lower'): OpDef(Label('lower'), [STR_TYPE], STR_TYPE, op_lower),
    Label('upper'): OpDef(Label('upper'), [STR_TYPE], STR_TYPE, op_upper),
    Label('ieq'): OpDef(Label('ieq'), [INT_TYPE, INT_TYPE], BOOL_TYPE, op_ieq),
    Label('ine'): OpDef(Label('ine'), [INT_TYPE, INT_TYPE], BOOL_TYPE, op_ine),
    Label('ilt'): OpDef(Label('ilt'), [INT_TYPE, INT_TYPE], BOOL_TYPE, op_ilt),
    Label('ile'): OpDef(Label('ile'), [INT_TYPE, INT_TYPE], BOOL_TYPE, op_ile),
    Label('igt'): OpDef(Label('igt'), [INT_TYPE, INT_TYPE], BOOL_TYPE, op_igt),
    Label('ige'): OpDef(Label('ige'), [INT_TYPE, INT_TYPE], BOOL_TYPE, op_ige),
    Label('seq'): OpDef(Label('seq'), [INT_TYPE, INT_TYPE], BOOL_TYPE, op_seq),
    Label('sne'): OpDef(Label('sne'), [INT_TYPE, INT_TYPE], BOOL_TYPE, op_sne),
    Label('slt'): OpDef(Label('slt'), [INT_TYPE, INT_TYPE], BOOL_TYPE, op_slt),
    Label('sle'): OpDef(Label('sle'), [INT_TYPE, INT_TYPE], BOOL_TYPE, op_sle),
    Label('sgt'): OpDef(Label('sgt'), [INT_TYPE, INT_TYPE], BOOL_TYPE, op_sgt),
    Label('sge'): OpDef(Label('sge'), [INT_TYPE, INT_TYPE], BOOL_TYPE, op_sge)}
#
# end of op.py
# Id: op.py 4838 2021-11-14 12:02:46Z \BlindText{wke@IPM.EDU.MO}
\end{dispPy}

\section{State transition rules}

\PyFile{st.py}
\begin{dispPy}
#    (c) 2021 \BlindText{Wei Ke and Ka-Hou Chan}
#    license:          GPL-3
#
from bidict import dict_union, bidict_union

from subtype import Value, NULL_TYPE, VALUE_TYPES, subtype

from graph import (
    Mismatch,   Label,      gc_state,   push_state, pop_state,
    add_object_to_state,    add_value_to_state,     swing_state,
    find_var,   find_lvar,  find_attr,  find_lattr, extract_pattern,
    cons_match)

from pattern import PatternConj, PatternDisj, MatchStmt

from asx import (
    AssignStmt, IfStmt,     WhileStmt,  PrintStmt,  BlockStmt,
    VarDecl,    VarEnd,     VarExpr,    AttrExpr,   OpExpr,
    NewExpr,    AndExpr,    OrExpr)

from op import invoke_op

SCOPE_LABEL = Label('$')

def eval_value(x, sg):
    return add_value_to_state(sg, x)

def eval_var(x, sg):
    return find_var(sg, SCOPE_LABEL, x.label)

def eval_lvar(x, sg):
    return find_lvar(sg, SCOPE_LABEL, x.label)

def eval_attr(x, sg):
    p = eval_expr(x.expr, sg)
    return find_attr(sg, p, x.label)

def eval_lattr(x, sg):
    p = eval_expr(x.expr, sg)
    return find_lattr(sg, p, x.label)

def eval_op(x, sg):
    la, args = x
    return invoke_op(sg, la, [eval_expr(arg, sg) for arg in args])

def eval_new(x, sg):
    return add_object_to_state(sg, x.cla)

def eval_and(x, sg):
    left, right = x
    p = eval_expr(left, sg)

    g, tm, vm = sg
    if vm[p].value != True:
        return p

    return eval_expr(right, sg)

def eval_or(x, sg):
    left, right = x
    p = eval_expr(left, sg)

    g, tm, vm = sg
    if vm[p].value != False:
        return p

    return eval_expr(right, sg)

EXPR_TAB = {
    Value: eval_value,
    VarExpr: eval_var,
    AttrExpr: eval_attr,
    OpExpr: eval_op,
    NewExpr: eval_new,
    AndExpr: eval_and,
    OrExpr: eval_or}

def eval_expr(x, sg):
    return EXPR_TAB[type(x)](x, sg)

LEXPR_TAB = {
    VarExpr: eval_lvar,
    AttrExpr: eval_lattr}

def eval_lexpr(lx, sg):
    return LEXPR_TAB[type(lx)](lx, sg)

def st_var_decl(vd, sg):
    la, cla = vd
    sg = push_state(sg, SCOPE_LABEL)
    q = add_object_to_state(sg, NULL_TYPE)
    swing_state(sg, sg.layout.root, la, q)
    return gc_state(sg)

def st_var_end(ve, sg):
    la = ve.label
    return pop_state(sg, SCOPE_LABEL)

def st_print(pr, sg):
    def to_str(p):
        cla = sg.types[p]

        if cla is NULL_TYPE:
            return 'null'

        if cla in VALUE_TYPES:
            return str(sg.values[p].value)

        return '{}@({})'.format(cla.tag, p.id)

    rs = [to_str(eval_expr(x, sg)) for x in pr.args]
    print(', '.join(rs))
    return sg

def st_assign(s, sg):
    lx, x = s
    (p, la) = eval_lexpr(lx, sg)
    q = eval_expr(x, sg)
    swing_state(sg, p, la, q)
    return gc_state(sg)

def st_if(s, sg):
    x, thens, elses = s
    p = eval_expr(x, sg)
    if sg.values[p].value == True:
        return st_stmt(thens, sg)
    else:
        return st_stmt(elses, sg)

def st_while(s, sg):
    x, ws = s
    p = eval_expr(x, sg)
    while sg.values[p].value == True:
        sg = st_stmt(ws, sg)
        p = eval_expr(x, sg)

    return gc_state(sg)

def st_block(blk, sg):
    for s in blk.stmts:
        sg = st_scope(s, sg)

    return gc_state(sg)

def st_scope(s, sg):
    if type(s) is VarDecl:
        return st_var_decl(s, sg)

    if type(s) is VarEnd:
        return st_var_end(s, sg)

    return st_stmt(s, sg)

def st_match(s, sg):
    x, cas = s
    p = eval_expr(x, sg)
    pg = extract_pattern(sg, p)

    for junc, s, extra in cas:
        m = match_junc(pg, junc, extra.get())
        if m is not None:
            sg = push_state(sg, SCOPE_LABEL)
            for la, q in m:
                swing_state(sg, sg.layout.root, la, q)
            sg = st_stmt(s, sg)
            return pop_state(sg, SCOPE_LABEL)

    return sg

def match_junc(pg, junc, extra):
    if type(junc) is PatternConj:
        pgs1, fs1, rms1 = extra
        return match_conj(pg, pgs1, rms1)

    if type(junc) is PatternDisj:
        pgs1, fs1, rms1 = extra
        return match_disj(pg, pgs1, fs1, rms1)

    pg1, rm1 = extra
    return match_one(pg, pg1, rm1)

def match_patterns(pg1, pg2):
    g1, ts1 = pg1
    g2, ts2 = pg2

    try:
        f = cons_match(g2, g1, subtype, ts2, ts1)
        return f
    except Mismatch:
        return None

def match_one(pg, pg1, rm1):
    f = match_patterns(pg, pg1)

    if f is None:
        return None

    return [(la, f[q]) for la, q in rm1.items()]

def match_disj(pg, pgs1, fs1, rms1):
    for pg1 in pgs1:
        f = match_patterns(pg, pg1)
        if f is not None:
            rm1 = dict_union(rms1)
            f1 = bidict_union(fs1)
            return [(la, f[(set(f1.inv[f1[u]]) & set(f)).pop()]) for la, u in rm1.items()]

    return None

def match_conj(pg, pgs1, rms1):
    fs = [match_patterns(pg, pg1) for pg1 in pgs1]
    if any(f is None for f in fs):
        return None

    rm1 = dict_union(rms1)
    f = bidict_union(fs)
    return [(la, f[u]) for la, u in rm1.items()]

STMT_TAB = {
    MatchStmt: st_match,
    AssignStmt: st_assign,
    IfStmt: st_if,
    WhileStmt: st_while,
    PrintStmt: st_print,
    BlockStmt: st_block}

def st_stmt(s, sg):
    return STMT_TAB[type(s)](s, sg)

def st_program(prog, sg):
    return st_block(prog.block, sg)
#
# end of st.py
# Id: st.py 4841 2021-11-14 17:09:10Z \BlindText{wke@IPM.EDU.MO}
\end{dispPy}

\section{Types and subtyping rules}

\PyFile{subtype.py}
\begin{dispPy}
#    (c) 2021 \BlindText{Wei Ke and Ka-Hou Chan}
#    license:          GPL-3
#
from collections import namedtuple

from pp import PrfxPP

class SubtypeError(Exception): pass
class DuplicateClassError(SubtypeError): pass
class UndefinedClassError(SubtypeError): pass
class AttrTypeConflictError(SubtypeError): pass
class MinTypeError(SubtypeError): pass

Tag = namedtuple('Tag', ['id'])
Lazy = namedtuple('Lazy', ['tag'])

def is_anon_tag(tag):
    return not tag or tag.id[0] == '*'

CLA_TAB = {}

class Cla:
    def __init__(self, tag, supers = [], attrs = [], tags = []):
        if tag:
            if tag in CLA_TAB:
                raise DuplicateClassError()
            CLA_TAB[tag] = self

        self.tag = tag
        self.tags = set(tags)
        if not is_anon_tag(tag):
            self.tags.add(tag)

        self.attrs = dict(attrs)

        for su in supers:
            self.tags.update(su.tags)
            for la, ty in su.attrs.items():
                if la in self.attrs:
                    if self.attrs[la] != ty:
                        raise AttrTypeConflictError()
                else:
                    self.attrs[la] = ty

        self.hashcode_cache = self.compute_hash()

    def compute_hash(self):
        s = list(self.tags)
        s.sort()
        return hash(('Cla', tuple(s)))

    def resolve_lazy(self):
        lzs = [(la, ty) for la, ty in self.attrs.items() if type(ty) is Lazy]
        for la, ty in lzs:
            self.attrs[la] = CLA_TAB[ty.tag]
        return self

    def __eq__(self, y):
        if self is y:
            return True

        if not is_anon_tag(self.tag) and self.tag == y.tag:
            return True

        return self.tags == y.tags

    def __hash__(self):
        return self.hashcode_cache

    def __le__(self, y):
        if self is y:
            return True

        if not is_anon_tag(self.tag) and self.tag == y.tag:
            return True

        return self.tags >= y.tags

    def to_pp(self):
        pp = {}
        pp['class tag'] = None if self.tag is None else self.tag.id
        pp['supers'] = sorted(tag.id for tag in self.tags)
        pp['attrs'] = {la.id: cla.tag.id for la, cla in self.attrs.items()}
        return PrfxPP('Cla', pp)

    @staticmethod
    def get(tag):
        if tag not in CLA_TAB:
            raise UndefinedClassError()
        return CLA_TAB[tag]

    @staticmethod
    def reset():
        CLA_TAB = {}

    @staticmethod
    def inter(cs):
        return Cla(None, cs)

    @staticmethod
    def union(cs):
        tags = set.intersection(*(cla.tags for cla in cs))
        attrs = set.intersection(*(set(cla.attrs.items()) for cla in cs))
        return Cla(None, [], attrs, tags)

NO_TYPE = Cla(Tag('*TOP'))
NULL_TYPE = Cla(Tag('NULL'), set(), {})
INT_TYPE = Cla(Tag('INT'), set(), {})
STR_TYPE = Cla(Tag('STR'), set(), {})
BOOL_TYPE = Cla(Tag('BOOL'), set(), {})

VALUE_TYPES = {INT_TYPE, STR_TYPE, BOOL_TYPE}

Value = namedtuple('Value', ['cla', 'value'])

class ValueSet:
    def __init__(self, vs):
        self.values = set(vs)
        s = list(self.values)
        s.sort()
        self.vector = tuple(s)

    def __hash__(self):
        return hash(('ValueSet', self.vector))

    def __eq__(self, y):
        return self.vector == y.vector

def subtype(x, y):
    if x is NULL_TYPE:
        return True

    if type(x) is Cla and type(y) is Cla:
        return x <= y

    if type(x) is ValueSet and type(y) is ValueSet:
        return x.values <= y.values

    if type(x) is ValueSet and type(y) is Cla:
        return all(v.cla == y for v in x.values)

    return False

def min_type(x, y):
    if subtype(x, y):
        return x

    if subtype(y, x):
        return y

    raise MinTypeError()

def classof(x, path = []):
    if type(x) is ValueSet:
        if path:
            return None
        clas = set(v.cla for v in x.values)
        if len(clas) != 1:
            return None
        return clas.pop()

    if type(x) is Cla:
        cla = x
        for na in path:
            if na not in cla.attrs:
                return None
            cla = cla.attrs[na]
        return cla

    return None

def exists_ty_le_all(ts):
    clas = [classof(t) for t in ts]
    if any(cla is None for cla in clas):
        return False

    if any(cla in VALUE_TYPES for cla in clas):
        return len(set(clas)) == 1

    return True

def ty_inf(ts):
    clas = [classof(t) for t in ts]
    if any(cla in VALUE_TYPES for cla in clas):
        vss = [t for t in ts if type(t) is ValueSet]
        if vss:
            return ValueSet(set().union(*[vs.values for vs in vss]))
        else:
            return clas[0]
    else:
        return Cla.inter(clas)

def ty_sup(ts):
    clas = [classof(t) for t in ts]
    if any(cla in VALUE_TYPES for cla in clas):
        vss = [t for t in ts if type(t) is ValueSet]
        if vss:
            return ValueSet(set().intersection(*[vs.values for vs in vss]))
        else:
            return clas[0]
    else:
        return Cla.union(clas)
#
# end of subtype.py
# Id: subtype.py 4838 2021-11-14 12:02:46Z \BlindText{wke@IPM.EDU.MO}
\end{dispPy}

\section{Type-checking rules}

\PyFile{tc.py}
\begin{dispPy}
#    (c) 2021 \BlindText{Wei Ke and Ka-Hou Chan}
#    license:          GPL-3
#
from subtype import Cla, Value, subtype, min_type, classof, BOOL_TYPE, VALUE_TYPES

from op import is_op, get_op

from asx import (
    PrintStmt,  AssignStmt, IfStmt,     WhileStmt,  VarDecl,
    VarEnd,     VarExpr,    AttrExpr,   OpExpr,     NewExpr,
    AndExpr,    OrExpr,     Program)

from pattern import PatternConj, PatternDisj, MatchStmt

from graph import (
    PatternGraph, cons_pattern_graph, cons_match_conj,
    cons_match_disj, unzip4)

class Env:
    def __init__(self, outer = None, items = []):
        self.tab = dict(items)
        self.outer = outer

    def __contains__(self, k):
        return k in self.tab or (self.outer is not None and k in self.outer)

    def __getitem__(self, k):
        if k in self.tab:
            return self.tab[k]
        else:
            return self.outer[k]

    def __setitem__(self, k, v):
        self.tab[k] = v

    def top(self):
        return set(self.tab)

    def pop(self):
        return self.outer

    def show(self):
        print(self.tab)
        env = self.outer
        while env is not None:
            print(env.tab)
            env = env.outer

class TypeCheckingError(Exception): pass
class NodeTypeError(TypeCheckingError): pass
class NodeSubtypeError(TypeCheckingError): pass
class IncompatibleTypesError(TypeCheckingError): pass
class UndefVarError(TypeCheckingError): pass
class ParentObjectError(TypeCheckingError): pass
class AttrError(TypeCheckingError): pass
class OpError(TypeCheckingError): pass
class OpArgLenError(TypeCheckingError): pass
class OpArgTypeError(TypeCheckingError): pass
class VarEndError(TypeCheckingError): pass
class LeftExprError(TypeCheckingError): pass
class ClassError(TypeCheckingError): pass
class AssignTypeError(TypeCheckingError): pass
class CondTypeError(TypeCheckingError): pass
class ScopeError(TypeCheckingError): pass
class PrintArgTypeError(TypeCheckingError): pass
class BoolTypeError(TypeCheckingError): pass

def tc_node(p, env, es, types):
    if p in env:
        return env[p]

    if p not in types:
        raise NodeTypeError()

    t = types[p]
    env2 = Env(env, {p: t})

    for la in es.labels[p]:
        q = es.targets[(p, la)]
        s = tc_node(q, env2, es, types)
        if not subtype(s, classof(t, [la])):
            raise NodeSubtypeError()

    return t

def tc_graph(pg, env):
    (ns, es, p), types = pg
    t = tc_node(p, env, es, types)
    return t

def tc_ref(la, env, pg, refs):
    p = refs[la]
    (ns, es, p2), types = pg
    t = tc_node(p, env, es, types)
    return t

def tc_pattern(pattern, env):
    g, tm, rm = cons_pattern_graph(pattern)
    pg = PatternGraph(g, tm)
    t = tc_graph(pg, env)
    rtm = {la: tc_ref(la, env, pg, rm) for la in rm}
    return (t, rtm, pg, rm)

def tc_conj(patterns, env):
    ts, rtms, pgs, rms = unzip4([tc_pattern(pattern, env) for pattern in patterns])
    fs, tsd = cons_match_conj(pgs)
    rtm2 = {la: min_type(s, tsd[f[rm[la]]])
            for f, rtm, rm in zip(fs, rtms, rms) for la, s in rtm.items()}
    th = tsd[fs[0][pgs[0].layout.root]]
    pts = [min_type(t, th) for t in ts]
    return (th, pts, rtm2, pgs, fs, rms)

def tc_disj(patterns, env):
    ts, rtms, pgs, rms = unzip4([tc_pattern(pattern, env) for pattern in patterns])
    fs, tsc = cons_match_disj(pgs)
    rtm2 = {la: tsc[f[rm[la]]] for f, rtm, rm in zip(fs, rtms, rms) for la in rtm}
    th = tsc[fs[0][pgs[0].layout.root]]
    pts = [t for t in ts]
    return (th, pts, rtm2, pgs, fs, rms)

def tc_case(ca, env):
    junc, stmt, extra = ca

    if type(junc) is PatternConj:
        t, pts, rtm, pgs, fs, rms = tc_conj(junc.patterns, env)
        extra.put((pgs, fs, rms))
    elif type(junc) is PatternDisj:
        t, pts, rtm, pgs, fs, rms = tc_disj(junc.patterns, env)
        extra.put((pgs, fs, rms))
    else:
        t, rtm, pg, rm = tc_pattern(junc, env)
        extra.put((pg, rm))

    env2 = Env(env, rtm)

    tc_stmt(stmt, env2)
    return t

def tc_value(v, env):
    return v.cla

def tc_var(x, env):
    la = x.label

    if la not in env:
        print(la)
        env.show()
        raise UndefVarError()

    return env[la]

def tc_attr(x, env):
    y, la = x
    t = tc_expr(y, env)

    if type(t) is not Cla:
        raise ParentObjectError()

    if la not in t.attrs:
        print(la)
        raise AttrError()

    return t.attrs[la]

def tc_op(x, env):
    op, args = x
    if not is_op(op):
        raise OpError()

    d = get_op(op)
    if len(args) != len(d.par_types):
        raise OpArgLenError()

    for y, pty in zip(args, d.par_types):
        if tc_expr(y, env) != pty:
            raise OpArgTypeError()

    return d.res_type

def tc_new(x, env):
    if type(x.cla) is not Cla:
        raise ClassError()

    return x.cla

def tc_and_or(x, env):
    left, right = x

    if tc_expr(left, env) is not BOOL_TYPE or tc_expr(right, env) is not BOOL_TYPE:
        raise BoolTypeError()

    return BOOL_TYPE

EXPR_TAB = {
    Value: tc_value,
    VarExpr: tc_var,
    AttrExpr: tc_attr,
    OpExpr: tc_op,
    NewExpr: tc_new,
    AndExpr: tc_and_or,
    OrExpr: tc_and_or}

def tc_expr(x, env):
    return EXPR_TAB[type(x)](x, env)

def tc_match(m, env):
    expr, cas = m
    t2 = tc_expr(expr, env)

    for ca in cas:
        t = tc_case(ca, env)
        if not (subtype(t2, t) or subtype(t, t2)):
            raise IncompatibleTypesError()

    return env

def tc_var_decl(vd, env):
    la, cla = vd
    return Env(env, {la: cla})

def tc_var_end(ve, env, env_base):
    if env is env_base:
        raise VarEndError()

    if set([ve.label]) != env.top():
        raise VarEndError()

    return env.pop()

def tc_print(pr, env):
    for x in pr.args:
        tc_expr(x, env)

    return env

def tc_assign(assign, env):
    lx, x = assign

    if type(lx) not in {VarExpr, AttrExpr}:
        raise LeftExprError()

    t = tc_expr(lx, env)
    t2 = tc_expr(x, env)

    if not subtype(t2, t):
        raise AssignTypeError()

    return env

def tc_if(ifs, env):
    x, thens, elses = ifs

    if tc_expr(x, env) is not BOOL_TYPE:
        raise CondTypeError()

    tc_stmt(thens, env)
    tc_stmt(elses, env)
    return env

def tc_while(ws, env):
    x, s = ws

    if tc_expr(x, env) is not BOOL_TYPE:
        raise CondTypeError()

    tc_stmt(s, env)
    return env

def tc_stmt(s, env):
    if type(s) is MatchStmt:
        return tc_match(s, env)

    if type(s) is AssignStmt:
        return tc_assign(s, env)

    if type(s) is IfStmt:
        return tc_if(s, env)

    if type(s) is WhileStmt:
        return tc_while(s, env)

    if type(s) is PrintStmt:
        return tc_print(s, env)

    return tc_block(s, env)

def tc_block(blk, env):
    env2 = env
    for s in blk.stmts:
        env2 = tc_scope(s, env2, env)

    if env2 is not env:
        raise ScopeError()

    return env

def tc_scope(s, env, env_base):
    if type(s) is VarDecl:
        return tc_var_decl(s, env)

    if type(s) is VarEnd:
        return tc_var_end(s, env, env_base)

    return tc_stmt(s, env)

def tc_program(prog, env):
    return tc_block(prog.block, env)
#
# end of tc.py
# Id: tc.py 4841 2021-11-14 17:09:10Z \BlindText{wke@IPM.EDU.MO}
\end{dispPy}

\section{Test cases}

\PyFile{test\_cases.py}
\begin{dispPy}
#    (c) 2021 \BlindText{Wei Ke and Ka-Hou Chan}
#    license:          GPL-3
#
from pp import pprint

from subtype import (
    Tag,        Lazy,       Cla,        subtype,    INT_TYPE,
    STR_TYPE,   BOOL_TYPE,  NO_TYPE,    Value,      ValueSet)

from pattern import (
    LabeledPattern, ClassPattern,   PatternRef, PatternConj,
    PatternDisj,    MatchStmt,      Case,       Extra)

from graph import (
    Label,  cons_pattern_graph, layout_graph_to_pp,
    init_state_graph)

from asx import (
    VarDecl,    VarEnd,     PrintStmt,  AssignStmt, BlockStmt,
    IfStmt,     WhileStmt,  VarExpr,    AttrExpr,   NewExpr,
    OpExpr,     AndExpr,    OrExpr,     Program)

from tc import tc_stmt, tc_program, Env

from st import st_stmt, st_program

def test_subtype():
    print(
'''
----
---- subtype ----
----
''')

    Cla.reset()

    a1 = Label('a1')
    a2 = Label('a2')
    b1 = Label('b1')
    b2 = Label('b2')
    c1 = Label('c1')
    c2 = Label('c2')
    d1 = Label('d1')
    d2 = Label('d2')

    A = Cla(Tag('A'),
        [],
        {
            a1: INT_TYPE,
            a2: STR_TYPE
        })

    B = Cla(Tag('B'),
        [A],
        {
            b1: INT_TYPE,
            b2: STR_TYPE
        })

    C = Cla(Tag('C'),
        [A],
        {
            c1: INT_TYPE,
            c2: STR_TYPE
        })

    D = Cla(Tag('D'),
        [B, C],
        {
            d1: INT_TYPE,
            d2: STR_TYPE
        })

    pprint(D.to_pp())

    E = Cla.inter([B, C])
    pprint(E.to_pp())

    print(subtype(D, E))
    print(subtype(E, D))

    pprint(NO_TYPE.to_pp())
    print(subtype(A, NO_TYPE))

    ONE = Value(INT_TYPE, 1)
    TWO = Value(INT_TYPE, 2)
    THREE = Value(INT_TYPE, 3)
    TEN = Value(INT_TYPE, 10)
    HELLO = Value(STR_TYPE, 'Hello')

    VSET1 = ValueSet(set(Value(INT_TYPE, x) for x in range(10)))
    print(subtype(ValueSet(set([ONE, TWO])), VSET1))
    print(subtype(ValueSet(set([ONE, TWO, TEN])), VSET1))
    print(subtype(VSET1, INT_TYPE))

    F = Cla.union([B, C])
    pprint(F.to_pp())
    print(subtype(F, A))
    print(subtype(A, F))
    print(A == F)

    h1 = Label('h1')
    h2 = Label('h2')
    k = Label('k')

    H = Cla(Tag('H'),
        [],
        {
            h1: INT_TYPE,
            h2: A
        })

    K = Cla(Tag('K'),
        [],
        {
            k: H
        })

    o = Label('o')
    q = Label('q')
    x = Label('x')
    y = Label('y')

    pb = ClassPattern(
        K,
        {
            k : ClassPattern(
                H,
                {
                    h1: ClassPattern(INT_TYPE, {}),
                    h2: ClassPattern(B, {})
                })
        })

    pc = ClassPattern(
        K,
        {
            k : ClassPattern(
                H,
                {
                    h1: ClassPattern(INT_TYPE, {}),
                    h2: LabeledPattern(x, ClassPattern(C, {}))
                })
        })

    s1 = BlockStmt([
        VarDecl(o, K),
        VarDecl(q, D),
        AssignStmt(VarExpr(o), NewExpr(K)),
        AssignStmt(AttrExpr(VarExpr(o), k), NewExpr(H)),
        AssignStmt(AttrExpr(AttrExpr(VarExpr(o), k), h1), Value(INT_TYPE, 3)),
        AssignStmt(VarExpr(q), NewExpr(D)),
        AssignStmt(AttrExpr(VarExpr(q), a1), Value(INT_TYPE, 4)),
        AssignStmt(AttrExpr(VarExpr(q), b1), Value(INT_TYPE, 5)),
        AssignStmt(AttrExpr(VarExpr(q), c1), Value(INT_TYPE, 6)),
        AssignStmt(AttrExpr(AttrExpr(VarExpr(o), k), h2), VarExpr(q)),

        MatchStmt(VarExpr(o), [
            Case(PatternConj([pb, pc]), BlockStmt([
                VarDecl(y, E),
                AssignStmt(VarExpr(y), VarExpr(x)),
                PrintStmt([
                    AttrExpr(VarExpr(y), a1),
                    AttrExpr(VarExpr(y), b1),
                    AttrExpr(VarExpr(y), c1)
                    ]),
                VarEnd(y)
                ]), Extra())
            ]),

        MatchStmt(VarExpr(o), [
            Case(PatternDisj([pb, pc]), BlockStmt([
                VarDecl(y, F),
                AssignStmt(VarExpr(y), VarExpr(x)),
                PrintStmt([AttrExpr(VarExpr(y), a1)]),
                VarEnd(y)
                ]), Extra())
            ]),

        VarEnd(q),
        VarEnd(o)])

    env = tc_stmt(s1, Env())
    sg = st_stmt(s1, init_state_graph())


def test_fig2():
    print(
'''
----
---- figure 2 ----
----
''')

    Cla.reset()

    e = Label('e')
    l = Label('l')
    r = Label('r')

    w = Label('w')
    x = Label('x')
    y = Label('y')
    z = Label('z')

    T_lz = Lazy(Tag('T'))

    T = Cla(T_lz.tag,
        [],
        {
            e: INT_TYPE,
            l: T_lz,
            r: T_lz
        }).resolve_lazy()

    pprint(T.to_pp())

    pattern = ClassPattern(
        T,
        {
            e: ValueSet({Value(INT_TYPE, 0)}),
            l: LabeledPattern(w, ClassPattern(
                T,
                {
                    l: LabeledPattern(x, ClassPattern(T, {})),
                    r: LabeledPattern(y, ClassPattern(
                        T,
                        {
                            l: PatternRef(w),
                            r: PatternRef(z)
                        }))
                })),
            r: LabeledPattern(z, ClassPattern(T, {}))
        })

    g, tm, rm = cons_pattern_graph(pattern)
    pprint(layout_graph_to_pp(g))

    o = Label('o')

    s1 = BlockStmt([
        VarDecl(o, T),
        AssignStmt(VarExpr(o), NewExpr(T)),
        AssignStmt(AttrExpr(VarExpr(o), l), NewExpr(T)),
        AssignStmt(AttrExpr(VarExpr(o), r), NewExpr(T)),
        AssignStmt(AttrExpr(AttrExpr(VarExpr(o), l), l), NewExpr(T)),
        AssignStmt(AttrExpr(AttrExpr(VarExpr(o), l), r), NewExpr(T)),
        AssignStmt(AttrExpr(AttrExpr(AttrExpr(VarExpr(o), l), r), l), AttrExpr(VarExpr(o), l)),
        AssignStmt(AttrExpr(AttrExpr(AttrExpr(VarExpr(o), l), r), r), AttrExpr(VarExpr(o), r)),

        AssignStmt(AttrExpr(VarExpr(o), e), Value(INT_TYPE, 0)),
        AssignStmt(AttrExpr(AttrExpr(VarExpr(o), l), e), Value(INT_TYPE, 1)),
        AssignStmt(AttrExpr(AttrExpr(AttrExpr(VarExpr(o), l), l), e), Value(INT_TYPE, 2)),
        AssignStmt(AttrExpr(AttrExpr(AttrExpr(VarExpr(o), l), r), e), Value(INT_TYPE, 3)),
        AssignStmt(AttrExpr(AttrExpr(VarExpr(o), r), e), Value(INT_TYPE, 4)),

        MatchStmt(VarExpr(o), [
            Case(pattern, PrintStmt([
                AttrExpr(VarExpr(w), e),
                AttrExpr(VarExpr(x), e),
                AttrExpr(VarExpr(y), e),
                AttrExpr(VarExpr(z), e)
                ]), Extra())
            ]),
        VarEnd(o)])

    env = tc_stmt(s1, Env())
    sg = st_stmt(s1, init_state_graph())


def test_fig3():
    print(
'''
----
---- figure 3 ----
----
''')

    Cla.reset()

    a = Label('a')
    b = Label('b')
    c = Label('c')
    d = Label('d')
    e = Label('e')
    x = Label('x')
    y = Label('y')

    X = Cla(Tag('X'),
        [],
        {
            b: INT_TYPE,
            c: STR_TYPE,
            d: STR_TYPE
        })

    pprint(X.to_pp())

    Y = Cla(Tag('Y'),
        [X],
        {
            a: INT_TYPE,
        })

    pprint(Y.to_pp())

    Z = Cla(Tag('Z'),
        [X],
        {
            e: BOOL_TYPE,
        })

    pprint(Z.to_pp())

    W = Cla(Tag('W'),
        [Y, Z],
        {})

    pprint(W.to_pp())

    p1 = ClassPattern(
        Y,
        {
            a: ValueSet(set(Value(INT_TYPE, x) for x in range(3))),
            c: LabeledPattern(x, ClassPattern(STR_TYPE, {})),
            d: ClassPattern(STR_TYPE, {})
        })

    p2 = ClassPattern(
        Z,
        {
            c: ClassPattern(STR_TYPE, {}),
            d: ClassPattern(STR_TYPE, {}),
            e: ValueSet({Value(BOOL_TYPE, False)})
        })

    p3 = ClassPattern(
        X,
        {
            b: ValueSet({Value(INT_TYPE, 1)}),
            c: ValueSet({Value(STR_TYPE, 'apple')}),
            d: ValueSet({Value(STR_TYPE, 'banana')})
        })

    p4 = ClassPattern(
        Z,
        {
            b: ValueSet({Value(INT_TYPE, 1)}),
            c: LabeledPattern(y, ValueSet({Value(STR_TYPE, 'apple')})),
            d: PatternRef(y)
        })

    conj = PatternConj([p1, p2, p3])
    disj = PatternDisj([p3, p4])

    o = Label('o')
    q = Label('q')
    u = Label('u')
    v = Label('v')

    prog = Program(BlockStmt([
        VarDecl(o, W),
        VarDecl(q, STR_TYPE),
        VarDecl(u, BOOL_TYPE),
        VarDecl(v, INT_TYPE),
        AssignStmt(VarExpr(o), NewExpr(W)),
        AssignStmt(AttrExpr(VarExpr(o), a), Value(INT_TYPE, 1)),
        AssignStmt(AttrExpr(VarExpr(o), b), AttrExpr(VarExpr(o), a)),
        AssignStmt(AttrExpr(VarExpr(o), c), Value(STR_TYPE, 'apple')),
        AssignStmt(AttrExpr(VarExpr(o), d), Value(STR_TYPE, 'banana')),
        AssignStmt(AttrExpr(VarExpr(o), e), Value(BOOL_TYPE, False)),

        MatchStmt(VarExpr(o), [
            Case(p4, PrintStmt([Value(STR_TYPE, 'p4, This should not match.')]), Extra())
            ]),
        PrintStmt([Value(STR_TYPE, 'after p4'), VarExpr(q)]),
        MatchStmt(VarExpr(o), [
            Case(p4, PrintStmt([Value(STR_TYPE, 'p4, This should not match.')]), Extra()),
            Case(disj, BlockStmt([
                AssignStmt(VarExpr(q), VarExpr(y)),
                PrintStmt([Value(STR_TYPE, 'disj'), VarExpr(q)])]), Extra())
            ]),
        MatchStmt(VarExpr(o), [
            Case(conj, BlockStmt([
                AssignStmt(VarExpr(q), VarExpr(x)),
                PrintStmt([Value(STR_TYPE, 'conj'), VarExpr(q)])]), Extra()),
            Case(p1, BlockStmt([
                AssignStmt(VarExpr(q), VarExpr(x)),
                PrintStmt([Value(STR_TYPE, 'p1'), VarExpr(q)])]), Extra())
            ]),
        MatchStmt(VarExpr(o), [
            Case(p4, PrintStmt([Value(STR_TYPE, 'p4, This should not match.')]), Extra()),
            Case(p1, BlockStmt([
                AssignStmt(VarExpr(q), VarExpr(x)),
                PrintStmt([Value(STR_TYPE, 'p1'), VarExpr(q)])]), Extra())
            ]),
        MatchStmt(VarExpr(o), [
            Case(p4, PrintStmt([Value(STR_TYPE, 'p4, This should not match.')]), Extra()),
            Case(p2, BlockStmt([
                AssignStmt(VarExpr(v), AttrExpr(VarExpr(o), b)),
                PrintStmt([Value(STR_TYPE, 'p2'), VarExpr(v)])]), Extra())
            ]),
        MatchStmt(VarExpr(o), [
            Case(p4, PrintStmt([Value(STR_TYPE, 'p4, This should not match.')]), Extra()),
            Case(p3, BlockStmt([
                AssignStmt(VarExpr(u), AttrExpr(VarExpr(o), e)),
                PrintStmt([Value(STR_TYPE, 'p3'), VarExpr(u)])]), Extra())
            ]),
        VarEnd(v),
        VarEnd(u),
        VarEnd(q),
        VarEnd(o)]))

    env = tc_program(prog, Env())
    sg = st_program(prog, init_state_graph())


def test_gcd():
    print(
'''
----
---- gcd ----
----
''')

    Cla.reset()

    m = Label('m')
    n = Label('n')
    t = Label('t')

    s1 = BlockStmt([
        VarDecl(m, INT_TYPE),
        VarDecl(n, INT_TYPE),
        AssignStmt(VarExpr(m), Value(INT_TYPE, 210)),
        AssignStmt(VarExpr(n), Value(INT_TYPE, 120)),
        PrintStmt([Value(STR_TYPE, 'numbers'), VarExpr(m), VarExpr(n)]),

        IfStmt(
            OrExpr(
                OpExpr(Label('igt'), [VarExpr(m), Value(INT_TYPE, 200)]),
                OpExpr(Label('igt'), [VarExpr(n), Value(INT_TYPE, 200)])),
            PrintStmt([Value(STR_TYPE, 'one is > 200')]),
            PrintStmt([Value(STR_TYPE, 'none is > 200')])),

        IfStmt(
            AndExpr(
                OpExpr(Label('igt'), [VarExpr(m), Value(INT_TYPE, 200)]),
                OpExpr(Label('igt'), [VarExpr(n), Value(INT_TYPE, 200)])),
            PrintStmt([Value(STR_TYPE, 'both are > 200')]),
            PrintStmt([Value(STR_TYPE, 'one is <= 200')])),

        BlockStmt([
            WhileStmt(
                OpExpr(Label('ine'), [VarExpr(n), Value(INT_TYPE, 0)]),
                BlockStmt([
                    VarDecl(t, INT_TYPE),
                    AssignStmt(VarExpr(t), VarExpr(m)),
                    AssignStmt(VarExpr(m), VarExpr(n)),
                    AssignStmt(VarExpr(n), OpExpr(Label('mod'), [VarExpr(t), VarExpr(n)])),
                    VarEnd(t)])),
            PrintStmt([Value(STR_TYPE, 'gcd'), VarExpr(m)])
        ]),

        VarEnd(n),
        VarEnd(m)
        ])

    env = tc_stmt(s1, Env())
    sg = st_stmt(s1, init_state_graph())


if __name__ == '__main__':
    test_subtype()
    test_fig2()
    test_fig3()
    test_gcd()
#
# end of test\_cases.py
# Id: test\_cases.py 4841 2021-11-14 17:09:10Z \BlindText{wke@IPM.EDU.MO}
\end{dispPy}

\endinput
